% Options for packages loaded elsewhere
\PassOptionsToPackage{unicode}{hyperref}
\PassOptionsToPackage{hyphens}{url}
%
\documentclass[
  11pt,
]{article}
\usepackage{amsmath,amssymb}
\usepackage{iftex}
\ifPDFTeX
  \usepackage[T1]{fontenc}
  \usepackage[utf8]{inputenc}
  \usepackage{textcomp} % provide euro and other symbols
\else % if luatex or xetex
  \usepackage{unicode-math} % this also loads fontspec
  \defaultfontfeatures{Scale=MatchLowercase}
  \defaultfontfeatures[\rmfamily]{Ligatures=TeX,Scale=1}
\fi
\usepackage{lmodern}
\ifPDFTeX\else
  % xetex/luatex font selection
\fi
% Use upquote if available, for straight quotes in verbatim environments
\IfFileExists{upquote.sty}{\usepackage{upquote}}{}
\IfFileExists{microtype.sty}{% use microtype if available
  \usepackage[]{microtype}
  \UseMicrotypeSet[protrusion]{basicmath} % disable protrusion for tt fonts
}{}
\makeatletter
\@ifundefined{KOMAClassName}{% if non-KOMA class
  \IfFileExists{parskip.sty}{%
    \usepackage{parskip}
  }{% else
    \setlength{\parindent}{0pt}
    \setlength{\parskip}{6pt plus 2pt minus 1pt}}
}{% if KOMA class
  \KOMAoptions{parskip=half}}
\makeatother
\usepackage{xcolor}
\usepackage[margin=1in]{geometry}
\setlength{\emergencystretch}{3em} % prevent overfull lines
\providecommand{\tightlist}{%
  \setlength{\itemsep}{0pt}\setlength{\parskip}{0pt}}
\setcounter{secnumdepth}{5}
\usepackage{amsmath,amssymb}
\emergencystretch=3em
\setlength{\parskip}{0.4em}
\providecommand{\PSH}{\operatorname{PSH}}
\providecommand{\MA}{\operatorname{MA}}
\ifLuaTeX
  \usepackage{selnolig}  % disable illegal ligatures
\fi
\usepackage{bookmark}
\IfFileExists{xurl.sty}{\usepackage{xurl}}{} % add URL line breaks if available
\urlstyle{same}
\hypersetup{
  pdftitle={Relative tropical supports on curves, model atoms, and the unibranch finite-energy equation},
  pdfauthor={Primary-source reconstruction for P16.1},
  hidelinks,
  pdfcreator={LaTeX via pandoc}}

\title{Relative tropical supports on curves, model atoms, and the
unibranch finite-energy equation}
\author{Primary-source reconstruction for P16.1}
\date{5 October 2026}

\begin{document}
\maketitle

\section{Exact result and retained
task}\label{exact-result-and-retained-task}

The original P16.1 target is unchanged. It asks for the genuine
finite-energy equation on every reduced irreducible geometrically
unibranch projective variety over every algebraically closed trivially
valued characteristic-zero field, in every dimension, for every ample
line bundle. The proposed singular envelope proof has not been completed
in this reconstruction. Accordingly this paper is a \textbf{partial
result}, and its candidate flag is false. Ample real numerical Cartier
classes are included in the genuine global partial theorem below; the
new local algebraic constructions use rational polarization data.

The principal new analytic result proved here is a curve statement with
relative tropical contact supports. It does not require an ambient
convex extension. Together with a determinant calculation it proves
slicing-weight independence and the exact radial trace equality required
in Wu 6.12, whenever the preceding open-trace extension has been
obtained. The latter antecedent is kept visible. A second result
supplies a bounded genuine model test with positive Monge--Ampère mass
at any prescribed divisorial point. This removes the \emph{Dirac
equation} antecedent from a local domination argument, but cannot
replace its local mass, comparison and determination hypotheses.

We reconstruct the other explicit algebraic inputs and the global
operator separately. We also derive the calibrated facet sign in a
finite harmonic chart, and explain precisely what remains to apply it to
the whole stated local class. We do not claim that these inputs alone
prove higher-dimensional bounded local/global equality, Rossi's
principle, or singular envelopes.

Only the six primary sources listed in the input were read:
Boucksom--Jonsson 1801.08229v4 {[}BJ{]}, Wu 2609.27610v1 {[}W{]},
Chambert-Loir--Ducros 1204.6277v2 {[}CLD{]}, Gubler--Martin 1608.08030v3
{[}GM{]}, Wanner 1801.04713v3 {[}Wa{]}, and Thuillier tel-00010990v1
{[}Th{]}. Embedded generated arguments and reports are diagnostic data,
never premises. Source results below are used with their actual scope;
this is not an independent correctness or novelty verdict on any source.
All claims here are proof drafts, not exhaustive finite computations.
Hash and build observations are reported separately at the end.

\section{The actual general-variety global
operator}\label{the-actual-general-variety-global-operator}

\subsection{Models, energy, and measures before
envelopes}\label{models-energy-and-measures-before-envelopes}

Fix an algebraically closed trivially valued characteristic-zero field
\(k\), a reduced irreducible projective \(k\)-variety \(X\) of dimension
\(n\), and \(\omega\in\operatorname{Amp}(X)\subset N^1(X)_{\mathbb R}\).
Thus \(\omega\) may be \(c_1(L)\) for any ample line bundle. Put
\(V=(\omega^n)>0\). BJ's conventions, LF435--469, permit reduced
nonnormal varieties. Section 7 returns to this scope after the
smooth-only b-divisor subsection. Its operator is already available on
this singular domain, independently of geometric unibranchness or an
envelope property.

A test configuration is closed integrally in its \emph{prescribed
generic fiber}: on a chart \(A\) over \(k[t]\), use the integral closure
in \(A\otimes_{k[t]}k(t)\). Absolute normalization can change a
nonnormal generic fiber and is not substituted. Finiteness of this
relative closure, equivariance and preservation of the generic fiber are
BJ 1.10--1.12. The vertical Serre criterion 1.13 follows from the
determinant trick for the polar ideal: integral closedness forces every
associated polar point to have invertible maximal ideal, hence to be a
vertical DVR; conversely vertical depth two excludes poles at
higher-codimension associated points. Thus vertical components \(E\)
have well-defined \(b_E=\operatorname{ord}_E(t)>0\) and
\(v_E=b_E^{-1}\operatorname{ord}_E|_{k(X)}\).

On a common relatively integrally closed configuration, nef model
metrics represented by vertical Cartier divisors \(D_i\) have the
genuine mixed raw measure \[
 M^{\mathrm{raw}}=\sum_E b_E
       \bigl((\omega_{1,\mathcal X}+D_1)|_E\cdots
             (\omega_{n,\mathcal X}+D_n)|_E\bigr)\delta_{v_E}.             \tag{1}
\] Intersection uses fundamental cycles, including multiplicities. The
weights are nonnegative because the restricted classes are nef. The
central fiber cycle is rationally equivalent to a generic fiber, so its
total mass is \((\omega_1\cdots\omega_n)_X\). BJ 4.3 is the genuine
model nef/PSH criterion; it is not a smooth BFJ criterion applied to a
quotient.

Compactification over \(\mathbb P^1\) defines the symmetric pairing \[
 \langle(\theta_0,\phi_0),\ldots,(\theta_n,\phi_n)\rangle
       =\prod_i(\theta_{i,\overline{\mathcal X}}+D_i).
\] Projection formula gives independence of the dominating
configuration. Translation of entry \(i\) by \(c\) adds
\(c(\prod_{j\ne i}\theta_j)_X\). BJ 7.1 extends the pairing to ample psh
tuples by the infimum over dominating model majorants. Monotonicity is
model intersection positivity. For upper semicontinuity, choose one
majorant tuple with pairing below a specified threshold, and use
continuity of \(\sup(\phi_i-\psi_i)\) to bound nearby tuples by those
majorants plus small constants. The translation rule bounds the change.
This proves the extension and its decreasing-net continuity without
envelopes.

Define the actual trivial-reference energy and its finite-energy space
by \[
 E_\omega(\phi)=\frac{\langle(\omega,\phi)^{n+1}\rangle}{(n+1)V},\qquad
 \mathcal E^1(X,\omega)=\{\phi\in\PSH(X,\omega):E_\omega(\phi)>-\infty\}.
                                                               \tag{2}
\] Then \(E_\omega(0)=0\) and \(E_\omega(\phi+c)=E_\omega(\phi)+c\)
exactly. It is increasing, concave, usc and decreasing-continuous. BJ
7.8 proves finite mixed pairings on \(\mathcal E^1\) tuples: after
subtracting constants, compare the finitely many ample classes by
positive multiples. Its Lemmas 7.9--7.10 bound the energy of a sum in
terms of individual energies, and expansion of the sum bounds each mixed
pairing. The cone spans the space of class/potential differences, so BJ
7.14 extends the pairing multilinearly.

For \(h\in\operatorname{PL}(X)_{\mathbb R}\) prescribe \[
 \int h\,dM^{\mathrm{raw}}=
 \langle(0,h),(\theta_1,\phi_1),\ldots,(\theta_n,\phi_n)\rangle.          \tag{3}
\] In the nef-psh case this is a positive linear functional, bounded by
its mass times \(\|h\|_\infty\). Genuine model-PL density in
\(C^0(X^{\rm an})\) and Radon representation give its unique positive
Radon extension. Multilinearity gives the signed case. These are BJ
7.18--7.19, including agreement with (1), exact mass, mixed
integrability of finite-energy tests, and decreasing continuity. In
particular no operator is defined by assigning a desired inverse
measure.

Weighted decreasing continuity is BJ 7.21. The integral with test
\(\phi_0\) is the pairing \textbf{containing \(\phi_0\) minus the
pairing containing zero}. Both pairings converge. Applying the
increasing-lsc weighted-measure lemma to a translated \(-\phi_0\) proves
weak convergence of the weighted measures. This also corrects the
reversed signs in LF4119--4123. Quantitative mixed pairing estimates in
Section 7 give strong continuity. Positive finite-energy mixed measures
charge no pluripolar sets. The diagonal probability measure is
\[\MA_\omega(\phi)=V^{-1}(\omega+dd^c\phi)^n;\] its raw version has mass
\(V\), and \(\MA_\omega(0)=\delta_{v_{\rm triv}}\).

\subsection{Comparison, uniqueness, and the exact target
condition}\label{comparison-uniqueness-and-the-exact-target-condition}

At a strict winning model component, the maximum model agrees with the
winning model near that component. Their intersection weights agree.
Weighted decreasing continuity with test \(\max(\phi,\psi)-\psi\)
extends this to BJ 7.40: \[
 \mathbf1_{\{\phi>\psi\}}\MA(\max(\phi,\psi))
    =\mathbf1_{\{\phi>\psi\}}\MA(\phi).                                \tag{4}
\] Its limiting citation is weighted 7.21, rather than mere unweighted
continuity. Equal masses and (4), applied with \(\psi-\epsilon\), give
\[
 \MA(\psi)(\{\phi<\psi-\epsilon\})
    \le\MA(\phi)(\{\phi\le\psi-\epsilon\}).
\] Letting \(\epsilon\downarrow0\) gives the comparison inequality
\[\int_{\{\phi<\psi\}}\MA(\psi)\le\int_{\{\phi<\psi\}}\MA(\phi).       \tag{5}\]
The strict winning locus carries the measure of the winning potential.

Use \[I(\phi,\psi)=\int(\phi-\psi)(\MA(\psi)-\MA(\phi)),\qquad
 J_\psi(\phi)=E(\psi)-E(\phi)+\int(\phi-\psi)\MA(\psi).\] Mixed
intersection negativity gives \(I\ge0\) and
\(I/(n+1)\le J_\psi\le nI/(n+1)\). The energy difference formula is \[
 E(\phi)-E(\psi)=\frac1{(n+1)V}\sum_{j=0}^n
 \int(\phi-\psi)(\omega+dd^c\phi)^j\wedge(\omega+dd^c\psi)^{n-j}.       \tag{6}
\] BJ 10.3--10.4 prove that \(I=0\) implies a locally constant
difference. In our irreducible scope, normalize equal suprema and test
the separation estimate with \(\delta_v\) and \(\delta_{v_{\rm triv}}\).
Divisorial Diracs have finite dual energy by BJ 11.1. Equality at
divisorial points then gives global equality by BJ 4.22. Thus equal
measures give \textbf{one} additive real constant, before EP. Global
domination 10.6 also precedes EP: replace \(\phi\) by
\(\rho=\max(\phi,\psi)\in\mathcal E^1\). Under the a.e. inequality its
\(I\) with \(\psi\) is \(-\int(\rho-\psi)\MA(\rho)\le0\). Separation and
the a.e. equality remove the resulting constant. This only applies to
genuine global potentials.

For a positive Radon raw target \(\nu\) the exact required mass is
\(V\). For \(\mu=\nu/V\), define \[
 E^\vee_\omega(\mu)=\sup_{\psi\in\mathcal E^1(X,\omega)}
       \left(E_\omega(\psi)-\int\psi\,d\mu\right).                     \tag{7}
\] Finiteness of (7) is equivalent to integrability of every
\(\mathcal E^1\) potential, by BJ 9.20 in its irreducible scope. Here is
the operative proof with the signs fixed. Normalize \(\sup\psi=0\) and
put \(J=\sup\psi-E\). If integration were unbounded below on \(J\le C\),
choose \(\psi_j\) with \(\int\psi_j\,d\mu\le-2^j\). The decreasing
convex partial sums \(\sum_{j\le m}2^{-j}\psi_j+2^{-m}0\) have energy at
least \(-C\); their limit has finite energy and integral \(-\infty\), a
contradiction. The mixed positive expansion gives
\[I(t\psi,0)\le t[1-(1-t)^n]I(\psi,0)\le nt^2 I(\psi,0),\quad0\le t\le1.\]
Together with \(J\asymp I(\cdot,0)\), scaling by \(t=J(\psi)^{-1/2}\)
when \(J\ge1\) proves \(-\int\psi\,d\mu\le A\sqrt{J(\psi)}+B\).
Subtracting \(J\) proves (7). Conversely a finite (7) excludes an
infinite negative integral of any finite-energy potential. A solution
necessarily has finite (7): mixed integrability and the tangent
inequality from (6) show that it computes the supremum. Arbitrary
support is allowed throughout.

\subsection{The conditional existence theorem, with every antecedent
retained}\label{the-conditional-existence-theorem-with-every-antecedent-retained}

Assume the genuine envelope property (EP) for \(\omega\): the usc
supremum of each nonempty bounded-above family of global \(\omega\)-psh
functions is global \(\omega\)-psh. Then every target (7) has an
\(\mathcal E^1\) solution, unique modulo one constant. Conversely
surjectivity onto all such targets implies EP. This is the actual
general-variety BJ 12.8, not the smooth sufficient theorem 12.9.

EP supplies compactness by BJ 5.11. Bounded divisorial evaluations give
a Tychonoff subnet. Its usc tail suprema are psh by EP, and BJ's
nonnegligibility at divisorial points identifies their decreasing limit
with the limiting divisorial values. It is nontrivial and
sup-normalized. This proves compactness downstairs, without assuming
smoothness.

The replacement for smooth differentiation is BJ 8.5:
\[E^\downarrow(\phi+\epsilon f)=E(\phi)+\epsilon\int f\MA(\phi)
                           +O_f(\epsilon^2),\quad f\in\operatorname{PL}_{\mathbb R}, \tag{8}\]
uniformly over \(\phi\in\mathcal E^1\). Its proof decomposes \(f\) into
two auxiliary ample model potentials, uses model-volume estimates, and
expands mixed positive measures of fixed intersection mass. The volume
ratio changes by \(O(\epsilon)\) and its prefactor supplies another
\(\epsilon\). Decreasing approximation passes to finite energy. Apply it
to both signs. For continuous \(f\), uniform PL approximation with error
\(\eta\) changes energy by at most \(|t|\eta\). The derivative squeeze
has endpoints \(\int f_m\MA(\phi)\pm\eta\), not two minus signs. EP is
required to identify \(E^\downarrow(\phi+tf)\) with
\(E(P_\omega(\phi+tf))\) for unbounded \(\phi\). Thus the two-sided
envelope-energy derivative follows under its actual antecedent.

Take a sup-zero model maximizing sequence for (7). BJ 9.15 bounds \(J\)
by its dual deficit plus \(E^\vee(\mu)\). Compactness gives a weak
subnet; \(J\) lsc puts the limit in \(\mathcal E^1\). BJ 9.19 gives
\(L^1(\mu)\) convergence on bounded-\(J\) sets. Its proof approximates
the target first by MA of maximizing model potentials, using (8) and the
Legendre inequality before attainment. The integration error bound 9.17
is uniform on those sets, while model measures are finite sums of
divisorial evaluations. The correct identity for absolute differences is
\(|\phi_i-\phi|=2\max(\phi_i,\phi)-\phi_i-\phi\). Target integration
therefore converges and energy usc gives attainment. For a maximizing
net of deficit \(\delta_i\), (8) yields
\[\epsilon\int f(\MA(\phi_i)-\mu)\le\delta_i+C_f\epsilon^2.\] Divide by
\(\epsilon>0\), pass to the net limit, let \(\epsilon\downarrow0\), and
use \(-f\). A constant maximizing net at the attained potential gives
the genuine equation. The raw objective is \(VE-\int\phi\,d\nu\); its
translation invariance uses precisely raw mass \(V\).

BJ 10.14 gives target completeness before EP and 12.7 gives the
bi-Lipschitz MA embedding. Surjectivity makes it onto and gives strong
\(\mathcal E^1/\mathbb R\) completeness; adding supremum differences
gives full completeness. Conversely, bounded increasing finite maxima of
model competitors are Cauchy because (6) gives
\(J_{\phi_i}(\phi_j)\le n(E(\phi_j)-E(\phi_i))\) for \(i\le j\).
Completeness and divisorial determination identify the limit with the
envelope; the continuous-obstacle characterization gives EP. Hence this
paper's conditional theorem cannot satisfy the original existence task.
For \(n=0\), \(X=\operatorname{Spec}k\), \(V=1\), potentials are
constants and \(E(c)=c\); the equation is immediate.

\section{Corrected auxiliary local
conventions}\label{corrected-auxiliary-local-conventions}

For an integral affine algebra \(R\) over a complete field, write
\[\operatorname{AFS}(R)=\left\{\max\{\lambda,c_i\log|f_i|+C_i\}_i\right\},\]
where the listed \(f_i\) generate \(R\) as a ground-field algebra,
\(c_i>0\) are rational, and \(\lambda,C_i\) are rational. This is
precisely W 3.15(1). It has a finite floor, so it is finite and
continuous. All regularizable auxiliary statements here use
\textbf{upper growth} \(u\le h+O(1)\), not W 3.15(2)'s lower bound
\(h\le u+O(1)\). On \(\mathbb A^1\),
\(\max(\log|z|,-m)\downarrow\log|z|\) is an actual AFS sequence. No
finite floor can be below \(\log|z|(0)=-\infty\). This disproves the
literal lower-growth membership, not a corrected upper-growth theorem.

For a finite PL function on a relative tropical image, its finite affine
selections obey W 3.1(1): every summand is defined and finite on its
stated neighborhood. At an infinite coordinate its coefficient must be
zero. The corrected contact condition is
\[f(t')\ge\ell(t')+C\quad\hbox{near }t,\qquad f(t)=\ell(t)+C.              \tag{9}\]
The constant is the same in both expressions. The literal omission of
\(C\) in W 3.1(3) makes the condition vacuous for a continuous function
with a local lower bound. We use (9), explicitly as an auxiliary
correction; BJ's genuine global PSH and finite-energy problem are not
redefined.

Finite maxima and positive rational sums of AFS functions remain AFS.
Distribute finite maxima and clear logarithmic coefficients into regular
monomials. Keep each single-generator branch obtained by choosing floors
in all other factors in its original form. Its coefficient stays
positive rational and its intercept is rational, so the original
generator list is retained. The sum of floors is rational. Their
differences form a rational vector lattice since
\(\max(a-b,c-d)=\max(a+d,c+b)-(b+d)\). Constants belong to this
difference lattice: if \(h\) is AFS, \(h+c-h=c\) for rational \(c\).

\section{Relative tropical supports on smooth
curves}\label{relative-tropical-supports-on-smooth-curves}

\subsection{Lisse membership and the outgoing-slope
proof}\label{lisse-membership-and-the-outgoing-slope-proof}

\textbf{Curve theorem.} Let \(\Omega\) be complete, algebraically closed
and nontrivially valued, of residue characteristic zero, and let \(C\)
be a smooth algebraic \(\Omega\)-curve. A finite continuous function
\(b\) on an analytic open is locally subharmonic if it is a finite PL
selection in regular log-norm coordinates and satisfies (9) on the
relative tropical image. Consequently a decreasing sequence of these
functions on a common open has a subharmonic or identically \(-\infty\)
limit on each connected component.

Here is a proof without ambient convex smoothing. At a point where the
coordinates are finite, shrink so the finitely many relevant regular
functions are analytic units. At an infinite coordinate, restrict to a
finite affine chart containing the image point. Every selected affine
form on that chart has coefficient zero in that coordinate. This removes
that coordinate on a smaller neighborhood. Identically zero pulled
coordinates are removed in the same way, rather than assigned a finite
logarithm. All remaining coordinates are units.

For a strictly affinoid neighborhood inside this unit locus, choose one
strictly semistable model. Existence and its finite skeleton are Wa
2.2--2.3. Th 2.2.24 shows that the logarithm of a unit factors through
the retraction to this skeleton: the Shilov maximum inequality on each
retraction fiber, applied both to the unit and its inverse, is equality.
Its edge restriction is piecewise affine and its interior weighted
Laplacian is zero. Equivalently Th 2.3.20 makes each unit log harmonic.
All the finitely many units use this \emph{same} model; refinements
preserve their pullbacks. Compactness allows a finite common
subdivision.

The function \(b\) also factors through the skeleton. On a subdivided
edge, its coordinates are affine. A continuous selection of finitely
many affine functions on an interval can change selection only at
intersections of two distinct affine functions. There are finitely many
such intersections; coincident affine functions give the same value.
Subdivide there. This proves piecewise affineness, including real
breakpoints which need not be formal vertices. It is precisely Wa 2.7
lisse membership. Local lisse membership glues; finite compact covers
give the same property on smaller strict affinoids. At a type I point
after discarding zero coordinates, unit norms are constant on a small
disk. At a type IV point finitely many unit-log germs are constant by Th
2.3.21's proof, LF3604--3607. A continuous finite selection of constant
values is locally constant on a connected smaller neighborhood. These
cases do not create a skeleton coefficient.

At any interior skeleton point \(x\), pull back the relative support in
(9) to a harmonic function \(H=\ell(-\log|g_i|)+C\). Then \(b-H\) has
local minimum zero. Every outgoing one-sided edge slope satisfies
\(\partial_e(b-H)\ge0\). The unit-log balancing identity gives
\(\sum_e w_e\partial_e H=0\), with positive edge weights. Therefore
\[\Delta b(\{x\})=\sum_e w_e\partial_e b
                   =\sum_e w_e\partial_e(b-H)\ge0.                   \tag{10}\]
At a type III breakpoint there are two outgoing directions and the same
argument applies. At any other interior point the Laplacian is zero.
Affinoid boundary coefficients are not tested for subharmonicity on the
open set. Wa 2.11, or Th 3.4.4 with its maximum-principle proof, now
gives subharmonicity. The sheaf property gives the conclusion on the
open. Th 3.1.8(ii) gives the decreasing-limit statement, including the
separate identically \(-\infty\) alternative.

An algebraic map and a complete field extension pull back the regular
functions and preserve their norm values at the restricted analytic
point. They preserve the finite selections and the relative support
inequality on a smaller image. Thus this theorem applies to the curve
restrictions of the corrected local decreasing tropical-PL class. It
does not require that the relative support extend convexly to an ambient
Euclidean open.

For clarity about that distinction, let \(S\) be the three positive rays
generated by \((1,0),(0,1),(1,1)\). Let \(f\) be zero on the first two
and \(f(t,t)=t\) on the third. It has relative local supports, including
zero at the origin. A convex ambient extension would satisfy
\(f(t/2,t/2)\le(f(t,0)+f(0,t))/2=0\), a contradiction. This unbalanced
image is not asserted to be a tropicalization of a boundaryless variety;
it only prevents an invalid general ambient-extension inference.

Wa Section 4 additionally assumes a smooth \textbf{proper} algebraic
curve. When using Wa 4.3 to identify a lisse subharmonic function with a
positive CLD current, or Wa 4.7 in the converse direction, first
compactify the smooth curve by the normalization of its projective
function-field model. Over the perfect \(\Omega\) this regular curve is
smooth, and the original curve is an open subcurve. Functions are only
considered on the original open, without extension across its new
boundary. Thus Wa's exact scope is met. The slope proof itself uses
Section 2 and Th and does not require 4.7 as an extra antecedent.

\subsection{The radial finite-trace boundary
equality}\label{the-radial-finite-trace-boundary-equality}

Let \(F\) be any complete residue-characteristic-zero field,
\(a\in C(F)\) on a smooth curve, and \(q:D\to C\) a finite map from a
disjoint union of smooth curves, each dominating \(C\). Suppose \(v\)
and \(w\) near the center and its finite preimage are usc decreasing
limits of corrected relative support PL functions on common
neighborhoods. Identically \(-\infty\) restrictions are permitted and
treated separately. Assume that on a punctured algebraic neighborhood
\[v(t)=\sum_{x\mid t}\dim_{H(t)}(\mathcal O(D)\otimes H(t))_x\,w(x).     \tag{11}\]
Then (11) holds at \(a\).

For nontrivial \(F\), take its completed algebraic closure \(\Omega\).
For trivial \(F\), first extend to \(F((T))\) with \(|T|=e^{-1}\) and
then take the completed algebraic closure. These are complete
algebraically closed nontrivially valued extensions; residue
characteristic remains zero. The curve theorem applies after extension
and turns the pulled sequences into subharmonic functions or identically
\(-\infty\) restrictions. Artinian trace base change is proved
independently below, so (11) pulls back without an AFS-regularizability
antecedent. The geometric fiber now consists of rational points.

For a subharmonic function \(h\) at a rational center, let \(\xi_R\) be
the single boundary Gauss point of a sufficiently small closed
coordinate disk, \(R\in|\Omega^\times|\). Disk comparison with the
constant harmonic function gives \(h(0)\le h(\xi_R)\); Th 3.1.9
guarantees finite boundary values for a nontrivial subharmonic function.
Since \(\xi_R\to0\), usc gives the reverse upper limiting inequality.
Hence \(h(\xi_R)\to h(0)\), also if \(h(0)=-\infty\). This is a boundary
Gauss limit, not an ordinary punctured-neighborhood limit. The
identically \(-\infty\) case is immediate.

Choose target coordinate \(t\) and source coordinates \(z_i\) at the
finitely many centers. On small disks \(t=c_i z_i^{e_i}U_i(z_i)\),
\(U_i(0)=1\), where \(e_i\) is the source DVR order of \(q^*t\). Shrink
until \(|U_i-1|<1\). Residue characteristic zero implies that every
nonzero rational number has norm one. Thus the binomial series for
\(U_i^{1/e_i}\) converges; its coefficients are bounded. The change
\(y_i=z_iU_i^{1/e_i}\) has linear coefficient one and is invertible on a
smaller disk by a contraction for the higher terms. We obtain
\(t=c_i y_i^{e_i}\).

A target Gauss point of small radius \(r\in|\Omega^\times|\) has one
source Gauss point in each disk, of radius \(R_i=(r/|c_i|)^{1/e_i}\).
Rescale to unit radii. The residue equation \(Y^{e_i}=\alpha T\) is
Eisenstein at \(T\) over \(\widetilde\Omega(T)\), so its residue degree
is \(e_i\). The polynomial field extension already has degree \(e_i\);
consequently this is a single analytic fiber point of actual
fiber-algebra dimension \(e_i\). It is not an unweighted count. At the
center the Artinian dimension is likewise \(e_i\).

Finite properness excludes other preimages for small target disks: the
closed complement of chosen open source disks has closed image avoiding
the center. Shrink the target into its complement. Its Gauss points also
avoid the finite algebraic exceptional set in (11). Therefore
\[v(\xi_r)=\sum_i e_iw(\xi_{R_i}).\] The radial limits prove the center
equality. If one limit is \(-\infty\), usc bounds the finitely many
other terms above on small neighborhoods, so the sum tends to
\(-\infty\). No subtraction of infinities occurs. Aggregate geometric
Artinian dimensions and descend the equality to \(F\). This proves the
required boundary identity without arbitrary-\(H(s)\) affine
regularization.

\section{Determinants, base change, and fixed slicing
weights}\label{determinants-base-change-and-fixed-slicing-weights}

\subsection{Finite-flat traces are actual affine FS
functions}\label{finite-flat-traces-are-actual-affine-fs-functions}

Let \(A\) be an integral affine algebra over a complete
characteristic-zero field, and \(B\) a reduced finite locally free
\(A\)-algebra of rank \(d\). At an analytic base point \(s\), put
\(F=H(s)\). The actual fiber algebra is
\[B\otimes_A F=\prod_{y\mid s} B_{s,y},\qquad
 m_y=\dim_F B_{s,y},\qquad\sum_y m_y=d.\] Nilpotents and residue degrees
are retained. Its residue fields are finite complete valued extensions
of \(F\), with their unique extended norms. Define
\(N(u)(s)=\sum_y m_yu(y)\).

For \(u=\max(\lambda,c_i\log|f_i|+C_i)\in\operatorname{AFS}(B)\), choose
integer \(D>0\) clearing the \(c_i\). Let \(b_0=1\), \(b_i=f_i^{Dc_i}\),
\(w_0=D\lambda\), \(w_i=DC_i\). Multiplication by \(\sum_iT_i b_i\) on
the locally free module \(B\) induces a scalar endomorphism on
\(\det_A B\). Consequently its determinant is a globally defined
polynomial \[P(T)=\det(\operatorname{mult}_{\sum T_i b_i})
      =\sum_{|\beta|=d}a_\beta T^\beta\in A[T],\qquad a_{(d,0,\ldots,0)}=1. \tag{12}\]
A globally free basis is unnecessary.

After a complete algebraically closed scalar extension of \(F\), each
local Artinian factor has a residue-field composition series.
Multiplication by \(\sum T_i b_i\) acts on each quotient by its residue
scalar. Therefore
\[P_s(T)=\prod_x\left(\sum_iT_i b_i(x)\right)^{\ell_x}.\] Above an
original \(y\) the total geometric lengths are \(m_y\): its separable
residue field gives \([H(y):F]\) embeddings, each repeated by the
original Artinian length. This is exactly the required analytic fiber
dimension.

For every positive real polyradius \(r\), the weighted Gauss norm is
multiplicative. To prove this without a value-group restriction, choose
the lexicographically largest exponent attaining each polynomial's
maximum weighted norm. In the coefficient at their sum, there is exactly
one product of maximal terms: any other maximal exponent in either
factor is smaller, and their sum is then smaller. All remaining
contributions have strictly smaller weighted norm. They cannot cancel
that contribution. The reverse bound is the ultrametric triangle
inequality. Isometric scalar extension preserves coefficient norms term
by term. Applying this to (12) with \(r_i=e^{w_i}\) gives \[
 N(u)=\max_\beta\left\{D^{-1}\log|a_\beta|
                         +D^{-1}\sum_i\beta_iw_i\right\}.              \tag{13}
\] Zero coefficients are discarded. The \(T_0^d\) coefficient gives the
rational floor \(d\lambda\). All logarithmic coefficients are positive
rational.

The \(a_\beta\) need not generate \(A\), so (13) alone is not AFS
membership. Choose algebra generators \(g_j\) of \(A\). Their images in
\(B\) are polynomials in the source generators \(f_i\). Let
\(Q=\max(0,\log|f_i|)\) and \(c_* =\min_i c_i>0\). Floors and intercepts
give a rational \(M_0\) with \(u\ge c_*Q-M_0\). A polynomial expression
of degree \(e_j\) gives \(\log|g_j\circ\pi|\le e_jQ+L_j\), where \(L_j\)
bounds coefficient norms. Choose rational \(\eta_j>0\) with
\(\eta_je_j\le c_*\) and rational \(M_j\) large enough. Then
\(u\ge\eta_j\log|g_j\circ\pi|-M_j\) everywhere. Summing on the fiber
gives \[N(u)\ge d\eta_j\log|g_j|-dM_j.\] Adding these dominated branches
leaves (13) unchanged and supplies the original base generators. This
proves \textbf{actual} AFS membership.

Fixed finite fiber sums preserve supplied decreasing AFS sequences, also
when a summand tends to \(-\infty\). They have local upper bounds from
their first continuous term, so \(+\infty-\infty\) never occurs. Upper
growth descends along finite maps: for an integral source coordinate a
monic equation gives the root bound
\(\log|f|\le\max_{j<m}(\log|a_j|/(m-j))\); polynomial degree bounds on
its base coefficients give upper AFS growth. Generic finiteness of a
trace limit is a separate obligation: a fiber all of whose values are
finite is sufficient, while finiteness at just one source point is
insufficient. If the limit is finite continuous, Dini applies on compact
inverse images and the trace error is at most \(d\) times the source
error. Traces are then continuous. Dini is not applied to unbounded
singular limits.

For example, \(A=k[a]\), \(B=k[z]\), \(a=z^2\), and
\(u=\max(\log|z|,-1)\) give \(P=T_0^2-aT_1^2\) and
\(N(u)=\max(-2,\log|a|)\). At zero the one-point nilpotent fiber has
dimension two and contributes \(2u(0)=-2\). Split points or a degree-two
residue field give the same value away from zero. An unweighted point
count fails.

\subsection{Artinian trace base change without any regularization
antecedent}\label{artinian-trace-base-change-without-any-regularization-antecedent}

For a scalar-extended base point \(s'\) above \(s\), the fiber is
\((B\otimes_A H(s))\otimes_{H(s)}H(s')\). Original idempotents first
split this algebra by original \(y\); each factor then splits into
refined local Artinian factors. Their dimensions over \(H(s')\) sum to
the dimension of the original factor over \(H(s)\). Every refined norm
restricts to the original norm since finite residue-field norm extension
is unique. Hence \[\sum_{x\mapsto y}\dim_{H(s')}B'_{s',x}=m_y,\qquad
 N_{\pi'}(a\circ r)(s')=N_\pi(a)(s)                                    \tag{14}\]
for every function \(a\) on the original fiber, including \(-\infty\)
values with local upper bounds. This also holds for nonflat finite
fibers when their actual dimensions are used. It proves the base change
used in (11) without assuming any FS approximation of \(w\).

\subsection{Fixed nonflat weights from global FS
tests}\label{fixed-nonflat-weights-from-global-fs-tests}

Now retain the original algebraically closed trivial \(k\), take
integral affine \(Y\) of dimension \(n\ge1\), and a finite surjection
\(\pi:Y\to S\), where \(S\) is an affine open of \(\mathbb A_k^n\). Its
ring \(A\) is normal; put \(B=k[Y]\), \(K=\operatorname{Frac}A\),
\(L=\operatorname{Frac}B\) and \(d=[L:K]\). Choose a dense principal
open \(W\) on which the map is étale.

For \(s\in S^{\rm an}\) take its \(F=H(s)\)-rational lift \(s'\) and a
line through \(s'\) generically in \(W_F\). It exists since \(F\) is
infinite and a nonzero polynomial cannot vanish on every line direction
through a point. Intersect the line with \(S_F\). Its preimage is cut in
the geometrically integral \(Y_F\) by \(n-1\) equations. The height
bound and the affine dimension formula give every minimal component
dimension at least one; finiteness gives at most one. Thus no point is
an isolated minimal zero-dimensional component. Reduce and normalize
this curve preimage. The resulting \(D\to C\) is finite. Each component
dominates the smooth line, and its module over a target DVR is
torsion-free; hence it is flat. The generic étale fiber is unchanged and
the rank is \(d\). Normal curves over the perfect \(F\) are smooth.
Normalization covers every underlying point over \(s'\), giving positive
integers \(m_{y,C}\) at the original finite analytic fiber points,
totaling \(d\).

For a global AFS test \(h\) on \(Y\), replace determinant (12) by the
generic field norm \(N_{L/K}(\sum T_i b_i)\). Its coefficients extend
across the nonflat locus. Indeed all conjugates of the \(b_i\) are
integral over \(A[T]\), so their product is integral; it lies in
\(K[T]\). Normality of \(A[T]\) puts the norm polynomial in \(A[T]\).
The generator argument for (13) gives a global AFS function \(v_h\) on
\(S\), agreeing with the actual trace on \(W\). Pull \(v_h,h\) to the
normalized slicing curve. Their relative support PL sequences and (14)
meet the curve theorem, so (11) gives
\[v_h(s)=\sum_{y\mid s}m_{y,C}h(y).                                    \tag{15}\]
The norm polynomial does not depend on \(C\) or \(W\).

Global AFS tests determine individual weights on this finite fiber.
Their difference lattice contains rational constants and separates
points: different analytic points have different regular-function
seminorms. Truncate the distinguishing log norm at a sufficiently low
rational floor, and add very low generator branches dominated at the two
points. This is an actual AFS test separating them. The closure of the
rational lattice on a finite set contains real constants and real affine
rescalings. For each pair it interpolates two values; finite minima and
maxima then interpolate every prescribed finite-set function. Thus it is
all real functions on the fiber. Equality of test integrals in (15)
gives equality of each individual weight. Write the resulting fixed
weights \(m_y(\pi)\). They are independent of the line, the étale open,
and the input tested. At a nonflat point they are \textbf{normalized
slicing degrees}, not necessarily dimensions of the original nonflat
fiber algebra.

It follows unconditionally that the fixed trace of a supplied decreasing
global AFS sequence is decreasing global AFS-regularizable, provided the
trace limit is generically finite. Global FS tests themselves have a
continuous trace through the whole nonflat locus by the same norm
formula.

For a weak local input \(u\) with proper birational witness
\(\mu:\widetilde Y\to Y\), first shrink \(W\) so \(\mu\) is an
isomorphism above \(\pi^{-1}(W)\). The exceptional image is a proper
closed subset of \(Y\); its finite image in \(S\) is closed of smaller
dimension. Choose the line \emph{after} this shrink. Its generic lift
through \(\mu\) extends at every closed point of its normalized smooth
curve by properness and the DVR valuative criterion; separatedness glues
the extensions. The curve theorem applies to the pullback of
\(u\circ\mu\), also allowing identically \(-\infty\) restrictions. If a
locally subharmonic, for example globally regularizable, trace extension
\(v\) on the smooth base has already been constructed agreeing on \(W\),
(11) proves its value at \(s\) equals the fixed weighted trace of \(u\).
The weights cannot depend on \(u\) because (15) proved their
independence first. A line chosen before the witness need not lift. This
repair removes the arbitrary-\(H(s)\) regularization antecedent
specifically from the boundary equality, \textbf{not} from the preceding
open-trace extension or its smooth weak regularization input.

The retained domain really can produce such nondiscrete internal fields.
Consider the map
\[Y=\operatorname{Spec}k[a,b,y,z]/(b^2-a^3)\longrightarrow
S=\operatorname{Spec}k[a,y,z].\] It is finite flat with basis \(1,b\).
At the target point with kernel \((a)\) and monomial weights
\(v(y)=1,v(z)=\sqrt2\), its residue completion contains the nondiscrete
value group \(\mathbb Z+\sqrt2\mathbb Z\); the pigeonhole principle
gives nonzero elements arbitrarily close to zero. The line \(a=\tau\),
with \(y,z\) fixed at the rational lift, has normalized preimage
\(\tau=t^2,b=t^3\) over that completion. This is an affine chart of the
projective cusp curve times two projective lines, a geometrically
unibranch projective threefold with ample product polarization. The
radial theorem applies to the needed equality over this field, but no
arbitrary-field AFS theorem is inferred.

\section{Global entropy supports and
localization}\label{global-entropy-supports-and-localization}

\subsection{Uniform simplex approximation and the real logarithmic
support}\label{uniform-simplex-approximation-and-the-real-logarithmic-support}

For nonzero regular functions \(f_i\) including affine generators, let
\(x_i=\log|f_i|\in\mathbb R\cup\{-\infty\}\) and choose rational \(a>0\)
with \(\log(r+1)/a\le1\). Set \[\chi(x)=a^{-1}\log(1+\sum_i e^{ax_i}).\]
It lies between \(\max(0,x_i)\) and that maximum plus one. On the
probability simplex \(\Delta\), with \(x_0=0\) and
\(F(p)=a^{-1}\sum_{i=0}^r p_i\log p_i\), Gibbs' inequality gives
\[\chi(x)=\sup_{p\in\Delta}\{\sum_{i\ge1}p_ix_i-F(p)\}.                 \tag{16}\]
The elementary proof is \(\log z\le z-1\) applied to the probability
ratios. At \(x_i=-\infty\) restrict to the face \(p_i=0\).

Take rational triangulations of \(\Delta\) with arbitrarily small mesh,
respecting every face and retaining its vertices. The lower convex
envelope of their lifted vertex values lies above \(F\) by convexity and
at most \(\eta\) above it by uniform continuity and barycentric
interpolation in one small simplex. Its dual is the finite vertex
maximum. Therefore the corresponding maximum \(h\) obeys
\(\chi-\eta\le h\le\chi\) on \textbf{all} finite \(x\), without a
bounded-coordinate restriction. Round each intercept downward to a
rational value with error at most \(\eta\), retaining the zero vertex
intercepts. The error is at most \(2\eta\) globally. Faces give the same
inequalities at negative infinite coordinates. Clear rational slopes
into regular monomials; the vertex branches retain zero and each
original generator. These are actual AFS functions. With errors
\(e_m=4^{-m}\), \(h_m+2e_m\) decreases to \(\chi\) globally, since it is
at least \(\chi+e_m\) while the next term is at most
\(\chi+2e_{m+1}<\chi+e_m\).

At contact \(x\), let \(p_0=1/Z\), \(p_i=e^{ax_i}/Z\), and omit zero
coefficients. Its support is
\[\ell(x')=-a^{-1}\sum p_i\log p_i+\sum_{p_i>0}p_ix_i'.\] Its gap times
\(a\) is relative entropy \(\sum p_i\log(p_i/p'_i)\), with unique zero
at \(x\), including faces. For uniform separation use \(t_i=-x_i\),
\(C=-a^{-1}\sum p_i\log p_i\), and \(D(t')=\chi(-t')-C+\sum p_i t_i'\).
If \(D\le b\), then \(\min t_i'\ge-(b+C)/p_0\) and
\(\sum p_i t_i'\le b+C\). Thus all coordinates have a finite lower
bound; positive-support coordinates have finite upper bounds, and the
others lie in compact extended intervals \([-B,+\infty]\). The gap is
lsc there, with value \(+\infty\) at a positive-support pole. Its
sublevels are compact, and their intersection as \(b\downarrow0\) is the
single contact point. They eventually lie in any contact neighborhood,
proving the strict uniform gap outside.

On countable \(k\), enumerate nonzero regular functions and let
\(\rho_j\) be these finite barriers for the generators and first \(j\)
functions. For \(v=\max(0,\log|z_i|)\), polynomial degree bounds give
\(\rho_j\le C_j(v+1)\). Choose positive rational \(\gamma_j\) with
\(\gamma_j\rho_j\le2^{-j}(v+1)\). The series \(\rho=\sum\gamma_j\rho_j\)
converges locally uniformly and has upper growth. The upper partial sums
\[U_N=\sum_{j\le N}\gamma_j\rho_j+2^{-N}(v+1)\] decrease to it. Finite
sums admit uniform global AFS approximants \(U_N-e_N\le W_N\le U_N\),
\(e_N=4^{-N}\). Then \(W_N+2e_N\) decreases globally to \(\rho\).
Omitting a fixed term has the same proof.

For the real support \(s=C+\sum_i\alpha_i\log|f_i|\), \(\alpha_i>0\),
choose \(d_i>0\) with \(\log|f_i|\le d_iv\) globally. Let rational
\(q_{i,m}\uparrow\alpha_i\), rational
\(B_m\downarrow A=\sum d_i\alpha_i\) strictly from above, and rational
\(C_m\downarrow C\). Define \(b_m=B_m-\sum d_iq_{i,m}>0\) and
\[s_m=C_m+\sum_i q_{i,m}\log|f_i|+b_mv.\] Away from poles,
\[s_m-s=(C_m-C)+\sum_i(\alpha_i-q_{i,m})(d_iv-\log|f_i|)+(B_m-A)v\ge0,\]
\[s_m-s_{m+1}=(C_m-C_{m+1})+(B_m-B_{m+1})v
              +\sum_i(q_{i,m+1}-q_{i,m})(d_iv-\log|f_i|)\ge0.\] At poles
both untruncated values are \(-\infty\), so compare values directly. For
rational \(\epsilon_m\downarrow0\), \(H_m=\max(s_m,\epsilon_mv-m)\) is
actual AFS and decreases to \(s\). Expand \(v\) and clear powers; the
second branch supplies the rational floor and original single-generator
branches. Combining this sequence with the remaining barrier series
gives a global AFS sequence for
\(\rho'=\gamma_j\ell+\sum_{i\ne j}\gamma_i\rho_i\). The initial seminorm
topology supplies a finite-list tropical neighborhood inside any
prescribed open. Hence \(\rho'\) is finite at contact, agrees there with
\(\rho\), and is below \(\rho-\delta\) outside that open. This proves
the exact supporting construction without singular analytic-PL gluing,
an APW premise, or singular EP.

\subsection{Actual poles repair the localization
estimate}\label{actual-poles-repair-the-localization-estimate}

Let \(R\) be integral affine over trivial \(k\), \(f\ne0\), \(L\)
nonempty compact in \(Y_f^{\rm an}\), and \(x\) an \(R\)-polynomial-hull
point in \(Y_f^{\rm an}\). Choose \(M\) with \(|\log|f||\le M\) on
\(L\cup\{x\}\). Suppose \(u\in\operatorname{AFS}(R)\), \(h\in R_f\) and
\[a\log|h|+b\le\max(u,-c\log|f|)+O(1),\qquad a,c>0.\] Then
\[a\log|h(x)|+b\le\sup_L(a\log|h|+b)+2cM.                              \tag{17}\]

On the finite normalization \(\overline R\) in the common fraction
field, evaluate on every prime divisor ray \(\tau\operatorname{ord}_D\).
Regular functions from \(R\) have nonnegative orders, so their AFS
maximum is uniformly bounded above and below on the ray. Division by
\(\tau\) forces
\(a\max(-\operatorname{ord}_D h,0)\le c\operatorname{ord}_D f\) whenever
\(h\) has a pole; no pole is allowed if \(\operatorname{ord}_D f=0\).
These function-field valuations are points of \(Y_f^{\rm an}\) even when
their centers lie on \(f=0\).

For positive integers \(p,q\) with \(q/p>c/a\), \(z=h^pf^q\) has
nonnegative order at every divisor of
\(\operatorname{Spec}\overline R\). The normal codimension-one
intersection criterion gives \(z\in\overline R\), hence integrality over
\(R\). As \(z\in R_f\), the finite algebra \(B=R[z]\) satisfies
\(B_f=R_f\). One \(f^N\) kills \(B/R\). Thus \(z^\ell f^N\in R\) for
every \(\ell\ge1\). Apply the \emph{original} \(R\)-polynomial hull
inequality to this regular function, multiply logarithms by
\(a/(p\ell)\), and use the two bounds by \(M\). The error is
\(2a(q/p+N/(p\ell))M\). Let \(\ell\to\infty\) and \(q/p\downarrow c/a\).
If \(h(x)=0\) the result is immediate; if \(h\) vanishes on \(L\), the
same hull test forces \(h(x)=0\). This proves (17), including all
\(-\infty\) endpoints and nonnormal algebras. Finite maxima give its AFS
version.

The printed denominator bound fails for \(f=z^2\), numerator \(g=z\),
\(h=g/f=z^{-1}\), \(u=0\), \(a=2,c=1\): the growth maxima agree although
the claimed \(ar\le c\) fails. Here \(g\notin(f)\) yet vanishes
everywhere on \(Z(f)\). Actual pole orders repair this intermediate
claim. To deduce the full regularizable maximum principle one still
needs controlled approximants on \(Y_f\) with growth at most
\(2\max(u,-c\log|f|)+O(1)\). For a supplied such sequence, (17) gives
\(4cM\). Decreasing usc suprema on compact \(L\) converge to the
limiting supremum by nested closed superlevel sets; ordinary Dini for a
singular limit is unnecessary. Letting \(c\downarrow0\) proves the
maximum principle under that approximation antecedent. That antecedent
is not proved by (17) itself.

\section{Algebraic collars and their exact
limitation}\label{algebraic-collars-and-their-exact-limitation}

For the projective closure of an affine embedding, choose very ample
\(L\), a complete rational FS reference \(\theta\), and distinguished
section \(s_0\) with \(Y=X\setminus Z(s_0)\). Arrange the proper AFS
exhaustion \(g=-\log\|s_0\|_\theta\ge0\). An actual AFS expression
\(a=\max(\lambda,q_j\log|f_j|+A_j)\) has \(f_j=F_j/s_0^{d_j}\) with
global homogeneous sections \(F_j\); this follows by homogenizing its
polynomial expressions in the affine embedding. Fix rational \(c,C,R\),
\(c,R>0\), with \(a-cg\le C\) on the entire shell \(R\le g\le R+1\). Set
\(\psi=\max(a-cg,C)\) for \(g\le R\) and \(C\) outside.

Here is an unconditional representation in BJ's genuine \emph{rational}
model-PL lattice. Let \(G=\min(g,R+1)\) and
\(H_j=\max(\log\|F_j\|_{d_j\theta},-M_j)\). Clipped section norms are
differences of actual basepoint-free FS metrics: relative to the trivial
metric, the numerator maximum is
\(\max(\log\|F_j\|_{\rm triv},d_jp_\theta-M_j)\), and its second branch
contains the complete reference system. Subtract \(d_jp_\theta\). Tensor
powers clear rational data. The same applies to \(G\) using \(s_0\).
Choose rational \(M_j\) so that \(-q_jM_j+(q_jd_j-c)G+A_j\le C\) at both
endpoints of \([0,R+1]\). Then clipping cannot alter a branch above the
floor. Let \[A_*=\max\{C,\lambda-cG,\ q_jH_j+(q_jd_j-c)G+A_j\},\] choose
rational \(K_*>0\) with \(C+K_*\ge\sup_X A_*\), and set
\[F_*=\min\{A_*,C+K_*(R+1-G)\}.                                      \tag{18}\]
All signed coefficients are rational, as required by BJ 2.1's vector
lattice. On \(g\le R\) the cap is inactive; on the shell \(A_*=C\); for
\(g\ge R+1\), including the whole boundary divisor, \(G=R+1\) and the
cap is \(C\). Thus \(F_*=\psi\) everywhere. No arbitrary analytic PL is
declared to be equivariant algebraic model PL.

For these particular collars, the local metric weights are finite convex
maxima of algebraic unit logarithms wherever nonconstant, and a
semipositive-reference constant near the outer boundary. Convex
smoothing on finite coordinates gives positive CLD curvature there. At a
zero coordinate the floor removes its branch on a smaller neighborhood.
This proves current positivity for this explicit class, without the
general relative-image positivity gate below.

Represent the total metric by its equivariant configuration and extend
to \(k((t))\), \(|t|=e^{-1}\). BJ flag-ideal generators \(t^{-\ell}f\)
have norm \(e^\ell|f(p(x))|\) at \textbf{every} extended analytic point.
Thus the formal-model metric is the actual pulled-back metric, not only
its value on the Gauss section. Integral closure of an ideal does not
change these norms: a monic integral equation and the ultrametric root
bound control each added integral element. Relative closure keeps the
original nonnormal generic fiber. Clearing rational coefficients gives
an actual model bundle. CLD 18.8.6 preserves positive curvature under
this field extension.

To obtain a central curve degree inequality, take an analytic point
reducing to the generic point of the integral curve on the formal
special fiber. Formal reduction is surjective on this flat admissible
model. CLD 17.7.10 identifies the curve and the restricted bundle as a
residual pre-model. CLD 18.8.4 with one bundle gives nonnegative degree
there. It is applicable to integer PL metrics; no use of 18.8.8's
strictness ambiguity is needed in this direction. Field/component
extension degrees in projection formula are positive, so the same
inequality descends if a complete algebraic closure is used. Every
vertical integral curve has nonnegative degree; hence the represented
bundle is relatively nef and BJ 4.3 makes this particular collar
genuinely global psh. Nilpotent special-fiber structure is retained in
cycles and does not change the nef test on integral curves.

If a global AFS sequence \(a_i\downarrow u\) is already supplied and
\(u\le g+B\), then for \(c>1\) choose one outer shell with \(u-cg<C-1\).
The compact sets \(\{a_i-cg\ge C\}\) on the shell are nested with empty
intersection, so an entire tail satisfies the same collar inequality.
Formula (18) yields genuine global model psh truncations decreasing to
the desired bounded floor. Decrease floors to \(-\infty\); finiteness of
\(u\) somewhere ensures nontriviality. Global decreasing closure and BJ
12.18 apply. In a rational class, uniform approximation of nef models by
ample FS metrics with summable constant tails preserves decrease. Add a
complete very ample system with sufficiently low rational weights,
leaving the metric unchanged. Ratios by powers of \(s_0\) generate
\(k[Y]\) and its own section gives a rational floor. This proves
controlled AFS regularization in the legitimate
\textbf{already-regularizable} case. It does not establish the same
antecedent for arbitrary local or weak inputs.

For an already-Runge Hartogs domain containing the zero section, strict
polynomial polytope exhaustion gives \(G_j(y,w)=\sum_i f_{j,i}(y)w^i\)
and, on the zero-section domain, \[-\log d_{\Omega_r}(y)=\max_{j,i\ge1}
                         (\log|f_{j,i}(y)|-c_j)/i.                    \tag{19}\]
The Gauss supremum on the fiber disk proves both bounds for its radius.
Use strict inequality opens rather than the false general equality
\(\operatorname{int}\{|f|\le a\}=\{|f|<a\}\), which fails for a constant
constraint. Finite coefficients in (19) are global algebraic functions.
Adding \(\max(H_r,\epsilon_rg-r)\) supplies a rational floor and the
original generators. On each compact coordinate sublevel an entire tail
decreases; with upper growth, the common shell and (18) give decreasing
global truncations and actual AFS regularization. The nested exhaustions
and compact smaller closed fiber disks give convergence even for
infinite radius. This verifies the algebraic coefficient step when the
domain is already Runge. It does not prove the weak-input Hartogs domain
is Runge.

\section{Calibrated facets and the bounded-local
gate}\label{calibrated-facets-and-the-bounded-local-gate}

\subsection{What the harmonic-chart computation actually
proves}\label{what-the-harmonic-chart-computation-actually-proves}

CLD 7.8.4 proves vanishing of a boundary integral for a harmonic PL
chart; it is not itself a statement that all relative-support functions
are psh. Its proof, LF9594--9653, derives the facet balancing relation
from Green's formula 5.11.1. We can reconstruct the required sign and
positivity on a \emph{finite calibrated polyhedral chart}. This useful
restricted statement has the following precise assumptions. On a compact
domain inside the analytic open, a strongly supporting skeleton has a
finite cell decomposition adapted to the canonical calibration, to the
test, to all unit-log coordinates, and to the affine selections of
\(b\). The harmonic coordinates defining the test have been added to the
chart. On each top-dimensional cell \(\sigma\), \(b\) and the
coordinates are affine; on every interior facet \(\tau\) the chart and
test restrictions agree from all adjacent cells. Collapsed cells are
assessed in this refined chart.

Choose a positive facet-volume normalization and the transverse vector
\(e_{\sigma,\tau}\) pointing from \(\tau\) into \(\sigma\), scaled by
its positive canonical calibration. More explicitly, the calibration
weight can be absorbed into this transverse vector, as in CLD
LF9621--9638; if it is written separately it is a positive \(w_\sigma\).
For an affine harmonic \(H\), that calculation is
\[\sum_{\sigma\supset\tau}w_\sigma\partial_{e_{\sigma,\tau}}H=0.        \tag{20}\]
The vector is inward relative to its adjacent cell but is an outgoing
direction from the facet; in dimension one it is exactly the outgoing
direction used in (10). This convention is necessary for the sign.

Apply Green's formula cell by cell. The interior \(dd^cb\) is zero on
affine cells. The terms containing \(b\,d''\alpha\) cancel across a
facet by continuity and the harmonic-chart boundary balancing in CLD
7.8.4. The remaining derivative terms give the corner coefficient
\[\kappa_\tau(b)=\sum_{\sigma\supset\tau}
                      w_\sigma\partial_{e_{\sigma,\tau}}b_\sigma.     \tag{21}\]
The sign is verified by \(b(x)=|x|\) on a line: the two directions
contribute \(1+1\), giving positive Laplacian, not its negative. A
tangential component of a transverse vector changes both \(b\) and a
tangential contact support by the same amount; subtraction eliminates
it. Choose a relative support \(H\) at an interior point of \(\tau\).
Since \(b-H\) is affine on \(\tau\) and has a local minimum zero, its
tangential derivative is zero. Its derivative into every adjacent cell
is nonnegative by (9). Subtracting (20) gives
\[\kappa_\tau(b)=\sum_{\sigma\supset\tau}
                w_\sigma\partial_{e_{\sigma,\tau}}(b-H)_\sigma\ge0.    \tag{22}\]
The restriction of a strongly positive \((n-1,n-1)\) test to the facet
is a nonnegative calibrated density. Thus the sum of (21) times those
densities is nonnegative. Tests supported away from the analytic
boundary have no outer boundary terms. If the refined chart collapses a
facet to dimension below \(n-1\), its pulled-back facet-volume test
vanishes. If its rank is sufficient, use (20)--(22). One cannot discard
a cell just because the \emph{original} tropicalization had too small a
rank: refinement by the test coordinates is essential. This proves
positive curvature under the stated finite-chart assumptions.

The remaining application is substantive. W's corrected local convention
is a finite affine selection on relative images of arbitrary analytic
neighborhoods. To use the preceding calculation and model theorems for
every such selection, one must obtain finite rational analytic domains
where, after tensoring, it is a unit logarithm plus a constant, with a
compatible strongly supporting refined chart for every test. Max/min
formulas for globally algebraic inputs provide this explicitly; (18) is
one such class. A word such as ``PL'' does not supply the required
analytic domain covering for an arbitrary relative selection. This
general passage has not been fully proved in this paper. Thus (22) is
not promoted into an unconditional higher-dimensional local/global
theorem.

\subsection{Strictness and the narrower smooth
route}\label{strictness-and-the-narrower-smooth-route}

The actual source definitions are different. CLD 12.5.2 allows
\(\log|g|+c\), \(c\in\mathbb R\), for integer PL, but strict integer PL
is G-locally \(\log|g|\) on strict analytic domains. Proposition 17.8.4
explicitly assumes \textbf{strict integer PL}. Corollary 18.8.8 is
printed for integer PL and its proof cites 17.8.4 for (iii) to (i). The
displayed proof does not explain that change of hypothesis. This is an
application gap here, not a counterexample to the corollary. The
positive-current to residual-nef direction 18.8.4 and its
field-extension consequence 18.8.6 are stated for integer PL and can be
used without it. The separate 18.8.9 does impose nontrivial valuation
and is not substituted.

If a continuous rational analytic metric on a smooth connected
projective \(X/k\) already has a finite rational-domain representation
and positive curvature, there is a valid narrower route. Extend to
\(K=k((t))\), \(|t|=e^{-1}\). A common tensor power clears rational
slopes and rational constants; the latter are realized by powers of
\(t\). If the domains are strict, this is a strict PL metric in GM's
sense. GM 1.1 is now applicable: the field is nontrivially valued and
complete and the scheme is proper. It supplies an algebraic/formal model
metric. CLD 18.8.6 pulls positivity to \(K\), and the residual
central-curve argument above makes the model nef. On smooth \(X_K\),
this is genuine semipositivity in the discretely valued
characteristic-zero theory. BJ Remark 4.19 then returns the original
function to genuine global PSH. That remark assumes smoothness,
connectedness, characteristic zero and an ample class, all retained
here. It is not used on a singular quotient. GM is never applied
directly over the trivial ground field. General real local constants
cannot be cleared over \(k((t))\) this way; the rational hypothesis
matters.

For a continuous genuine global psh limit, decreasing genuine model
approximation is uniformly convergent by Dini. Model compatibility in
CLD 18.4.7 and continuous uniform convergence in 8.7.2 then identify its
mixed products with the global genuine products. CLD 18.4.7 retains
graded residue degrees, analytic multiplicities and defects; an
unweighted point count is not its formula. In contrast, CLD 8.7.2
explicitly assumes locally uniform convergence to \textbf{continuous}
functions. It does not justify a merely decreasing bounded discontinuous
limit.

\subsection{Correct common reference and what remains for mixed
products}\label{correct-common-reference-and-what-remains-for-mixed-products}

We verify the proposed collar bookkeeping independently. On an affine
chart choose finite rational coordinate inequalities, a complete ample
section system and a rational \textbf{negative} FS floor, giving a
proper local potential \(v\) with \(v(x)<-\epsilon\),
\(Q=\{v\le0\}\Subset U\) and \(P=\{v\le-\epsilon\}\). Section powers and
a common denominator express the finitely many algebraic constraints;
the complete system makes \(Q\) compact. Thresholds can be tightened
rationally using the initial seminorm topology. Including the
distinguished section with weight zero would force \(v\ge0\) and
invalidate this construction. It must have its negative weight. Let
\(\theta\) be this actual FS reference on the same rational ample class
and write \(\theta=\omega+dd^cp\) for its bounded global model-metric
difference from the requested reference.

Suppose \(|u|\le M\) on a neighborhood of \(Q\), with
common-neighborhood continuous corrected PL approximants
\(u_i\downarrow u\). Enlarge \(M\) to a rational bound. Choose rational
\(A\) with \(A\epsilon>2M+1\) and set \(h=Av+M+1\). On \(P\), \(h<-M\);
on a fixed compact collar \(v\ge-1/(2A)\), \(h\ge M+1/2\). The nested
compact superlevel sets of \(u_i-h\) on that collar have empty
intersection, so one whole tail is strictly below zero. The functions
\[\Psi_i=\max(u_i-h,0)\ \hbox{on }Q,\qquad \Psi_i=0\ \hbox{outside }Q\]
are continuous rational analytic PL and locally \(A\theta\)-psh. There
is no use of Dini to a discontinuous limit. If the finite-domain and
positive-curvature hypotheses of the preceding smooth route are proved
for these inputs, each \(\Psi_i\) is genuine global psh; decreasing
closure makes their bounded limit \(\Psi\) genuine global and finite
energy. Set \[\Phi=\Psi/A+p,\qquad G=v-p+(M+1)/A.\] On \(P\) the exact
identity is
\[\Phi+G=u/A.                                                        \tag{23}\]
\(\Phi\) is in the requested genuine \(\omega\) class. This retains both
the reference difference and the scalar. One common \(Q,P,A,p,G\) works
for a finite bounded mixed tuple and its bounded test by taking a common
bound and tail. Neither local MA nor comparison has been used to obtain
(23). However its genuine-global assertion for the whole local class
remains conditional on the just identified representation/positivity
passage.

For such \emph{supplied genuine common representations}, the intended
local product is \[M_{\rm loc}(u_1,\ldots,u_n)=A^n
       M_{\rm glob}((\omega,\Phi_1),\ldots,(\omega,\Phi_n))|_P.         \tag{24}\]
For continuous approximants, CLD/model agreement makes it intrinsic. On
fixed representation data, two different approximating sequences have
the same global limiting truncations \(\max(u_j-h,0)\); thus BJ
decreasing continuity gives the same limiting product. Different
representations can use one common input sequence on an overlap and the
continuous intrinsic product there, then pass to limits against
compactly supported tests. Weighted continuity uses
\(u_{0,i}=A(\Phi_{0,i}+G)\) and BJ 7.21; the fixed \(G\) is continuous
and bounded on compact chart subsets. This is the correct mechanism for
bounded discontinuous tests and their integrability. It is not an
invocation of CLD 8.7.2 for such limits. A common representation also
transfers the strict-max locality with the winning measure.

Still required for a global locally psh metric are the partition of
unity, integration by parts, total mass and comparison argument. W 4.14
starts by extending arbitrary local PL functions to global PL. Formula
(18) only does this for globally algebraic AFS inputs, whereas a local
test or analytic-piece bump need not have that form. In W 4.12 the
total-mass calculation moves \(dd^c\) across a PL partition and cancels
the resulting sum. Without justified global extensions, or a fully
constructed closed local current calculus and compatible weighted
convergence, this is not proved by positivity of the top-degree measures
in (24). No such missing step is silently assumed here. In particular
full smooth bounded local comparison/globalization, followed by the
smooth weak affine base, is not claimed as an unconditional new theorem
in this paper.

\section{Ample model atoms and the conditional smooth globalization
repair}\label{ample-model-atoms-and-the-conditional-smooth-globalization-repair}

\subsection{A genuine positive atom without a Dirac
equation}\label{a-genuine-positive-atom-without-a-dirac-equation}

\textbf{Model-atom proposition.} On every reduced irreducible projective
\(X/k\) with rational ample class \(c_1(L)\), every source-convention
divisorial point \(x\in X^{\rm div}\) has positive actual MA mass for
some bounded genuine ample model potential \(b\).

BJ 2.12 realizes \(x\) as \(v_E\) on the relative integral closure of a
flag-ideal blowup dominating \(X\times\mathbb A^1\). Its proof uses 1.20
and passage to a higher configuration, preserving that component's
valuation. BJ 2.33 supplies on this same model an ample test bundle:
choose a relatively ample vertical rational divisor \(D\) over the
trivial configuration and take \(\mathcal L=L_{\mathcal X}+\epsilon D\)
for small positive rational \(\epsilon\). The class is ample relative to
\(\mathbb A^1\), so its restriction to each central integral component
is ample. Hence \((\mathcal L|_E)^n>0\), and the weight
\(b_E(\mathcal L|_E)^n\) in (1) is strictly positive. Relative integral
closure retained the original generic fiber even when nonnormal.
Normalization by \(V\) preserves strict positivity. Multiple components
with the same valuation add positive weights. This proves the
proposition without EP or equation surjectivity. The source divisorial
convention uses rationally scaled valuations; no model atom at every
arbitrary real scaled point is asserted. The trivial valuation is
supplied by the trivial ample model. Dimension zero is the one-point
case already treated.

\subsection{What model atoms prove once local mass and comparison are
supplied}\label{what-model-atoms-prove-once-local-mass-and-comparison-are-supplied}

Suppose a bounded local \(\omega\)-psh cone on the projective space has
compatible mixed positive measures, exact raw mass \(V\), strict-max
locality, comparison (5), divisorial determination, and contains the
bounded genuine model potentials with their actual measures. These are
explicit antecedents. If bounded \(\phi\le\psi\) a.e. for the measure of
\(\psi\), fix divisorial \(x\), choose the model atom \(b\), and
translate it so \(b<\phi\) with a uniform strict margin. For \(0<t<1\)
put \(\phi_t=(1-t)\phi+tb\). Then \(\phi_t<\psi\) a.e. for the measure
of \(\psi\). Comparison and equal masses give
\[V=\MA^{\rm raw}(\psi)(\{\phi_t<\psi\})
   \le\MA^{\rm raw}(\phi_t)(\{\phi_t<\psi\})\le V.\] Mixed positivity
gives \(\MA^{\rm raw}(\phi_t)\ge t^n\MA^{\rm raw}(b)\), so its atom at
\(x\) is positive. The strict set has full measure, so
\(\phi_t(x)<\psi(x)\). Let \(t\downarrow0\) and use divisorial
determination to get \(\phi\le\psi\) everywhere. This rigorously
replaces W 4.16's continuous Dirac solution antecedent, but not its
preceding local theory.

On a smooth connected projective \(X\) in characteristic zero, known
smooth EP would then give bounded local/global equality. For any bounded
usc obstacle \(\phi\) on compact Hausdorff \(X^{\rm an}\) use the
directed net of all continuous majorants, ordered downward by finite
minima, not a sequence. Such majorants separate a point from closed
superlevel sets by normality: if \(c>\phi(x)\), choose an open
neighborhood on which \(\phi<c\) and a continuous cutoff equal to zero
at \(x\) and one off that neighborhood; interpolate between \(c\) and a
common upper bound. Finite minima give a decreasing net with infimum
\(\phi\). Common lower and upper truncations preserve these properties.

Smooth EP makes \(p_i=P_\omega(f_i)\) continuous genuine global psh.
Contact orthogonality follows from the actual derivative: the obstacle
\(p_i+t(f_i-p_i)\) has the same envelope \(p_i\) for \(0\le t\le1\), so
\(\int(f_i-p_i)\MA(p_i)=0\). Thus its measure is on the closed contact
set. On that set \(\phi\le f_i=p_i\). The conditional local domination
just proved gives \(\phi\le p_i\le f_i\). These decreasing genuine
envelopes limit to \(\phi\) with a common lower bound; genuine closure
gives global pshness. Bounded floors give the unbounded nontrivial case.
Nef references can be approached from above by ample ones using the
actual global class closure, provided the corresponding local hypotheses
hold. No sequence of majorants on an uncountable-field analytification
is asserted.

Likewise the smooth affine regularization has a noncircular
\textbf{conditional} route. On a smooth projective compactification with
effective ample boundary choose actual proper AFS \(g\). Boundary
blowups preserve an effective ample boundary divisor by taking a
sufficiently large multiple of its pullback minus the exceptional
divisor, whose negative is relatively ample. For a local input
\(u\le Ag+B\), choose rational \(A'>A\) and \(c>1\). The floors
\(\max(u-cA'g,C)\) extend constantly by \(C\) on one whole open boundary
collar since \(B-(cA'-A)g<C\) there. On the affine chart their local
reference weights are \(\max(u,C+cA'g)\), and near the boundary they are
semipositive-reference constants. The conditional bounded local/global
result makes them genuine. Decreasing floors and actual global FS
restriction give the sought affine sequence, with complete-system
generators and a floor. No already-regularizable extension lemma is used
to prove this antecedent.

For weak inputs one would resolve and compactify the witness to a smooth
projective birational map \(Z\to X\). The same floors upstairs lie in
the nef big pullback class; the smooth nef local/global statement is
needed there. BJ 5.13 is onto downstairs because the slightly larger
\emph{smooth} ample classes have EP. Analytic surjectivity identifies
the descended values. This chain is valid with the stated local
antecedents, which have not been supplied for the full class here.
Moreover regularization on \(Y_f\) alone does not yield the pole-free
controlled original-exhaustion approximants needed by the principal-open
removable extension. That is a separate gap, not covered by the model
atom or radial trace theorem.

\section{Rossi, the integer trace, and the remaining singular
induction}\label{rossi-the-integer-trace-and-the-remaining-singular-induction}

W 5.6 constructs on an arbitrary relatively compact analytic
neighborhood \(V\) a function equal to \(Cu\) outside and
\(\max(Cu,\log|f|-t)\) inside. Its boundary margin is correct, but its
next step invokes W 3.24 to approximate this analytic-piece PL function
globally by affine FS. Such a piece need not be a single global
algebraic maximum. Formula (18) supplies no extension on this arbitrary
analytic \(V\). We have neither an independent polynomial-hull proof of
Rossi's assertion nor a complete general analytic-piece model
representation here. The source's conclusion is not disproved by that
missing reconstruction.

There is a separate small logical error which can be repaired locally.
In W 6.15 LF2456--2458, continuity of a positive integer trace is used
to infer one integer \(\gamma\) on the entire chosen open. Continuity
gives \emph{local} constancy, not constancy on a possibly disconnected
open. An example occurs within the original field domain: take
\(z\mapsto a=z^2\) on \(\mathbb A^1_k\), with \(k\) algebraically closed
trivial characteristic zero. Let the base open be the disjoint union of
small disks about \(a=0\) and \(a=1\), with radius \(\rho<1\). Over the
first choose \(W_1\) to contain its whole preimage disk; over the second
choose \(W_1\) to contain the disk about \(z=1\) and \(W_2\) the disk
about \(z=-1\). The locally constant test is one on \(W_1\) and zero on
\(W_2\). Its actual finite-flat trace is two on the first base disk and
one on the second. At the branch center its Artinian length is two;
elsewhere all factors over the first disk are chosen, with dimensions
totaling two whether or not they split. Over the second the unit-root
expansion gives the two disjoint disks and trace one. Thus it is a
continuous integer trace on a disconnected open with no single
\(\gamma\). One may take smooth closed \(Z=\{0,1\}\) on the source; its
target images are distinct, and its other fiber point is \(-1\), so this
also fits the local split geometry.

The correct repair is to shrink around the tested base point to a
neighborhood on which the \emph{continuous} fixed trace is constant.
Divisorial determination, if already proved for that trace extension,
then identifies that integer at any remaining points of this smaller
neighborhood. A neighborhood crossing two different locally constant
values cannot be used. If \(\gamma=1\), the fiber subset in \(W_1\) has
total positive integer weight one, so it consists of the tested point
with weight one. The other weighted sum there is zero. In the negative
estimate of W Step 5, its trace is then at most \(-\delta\); no bound
\(\delta/(2(\gamma-1))\) is imposed and no division by zero is made. For
\(\gamma>1\), the original finite weighted bound works after this local
shrink. These repairs do not construct the arbitrary analytic-piece
function \(\widetilde h\) or \(\widetilde v\) that the norm is applied
to.

The dependency chain is therefore still broken before the singular EP
theorem. W 6.15 uses analytic split-piece gluing, the regularization of
its weak-input trace on the smooth target, Rossi and the localization
maximum principle. The radial theorem repairs the final equality in
6.12, but does not supply that smooth-target extension. W 6.14 then uses
6.15 for the weak Hartogs domain to be Runge; (19) only treats an
already Runge domain. W 6.16's geometric local projection would be
needed in the singular-stratum induction 6.3, and 6.5 uses that
induction for weak/global identification on affine charts. We do not
accept these as reconstructed merely because their geometric outline is
plausible. In particular no unreconstructed APW assertion, or the
absence of its separately supplied body, is used as a reason for
mathematical falsity.

The actual body W 7.1, LF2697--2747, asserts the desired geometrically
unibranch envelope property. Its countable-field proof descends the
envelope on a smooth resolution and then explicitly uses weak/global
6.5. The old circular 4.19 invocation and lower-growth inconsistency are
not repaired by quoting this later theorem. The present replacements
still leave the bounded-local and Runge gates just identified. We
therefore cannot adopt its singular conclusion as an established
antecedent for BJ 12.8.

\section{Normalization, conditional cardinality removal, and a genuine
obstruction}\label{normalization-conditional-cardinality-removal-and-a-genuine-obstruction}

\subsection{One-way identities and the separate onto
obligation}\label{one-way-identities-and-the-separate-onto-obligation}

For reduced irreducible geometrically unibranch \(X\), its finite
normalization \(q:X^\nu\to X\) is birational of degree one. The local
geometric branch condition says that each geometric fiber is a singleton
with purely inseparable residue extension; characteristic zero makes
that extension an equality. Thus it is a universal homeomorphism. Its
analytification is a homeomorphism as well: an analytic point is a norm
on the residue field at its support, and the unique support point
upstairs has that same residue field. Properness turns the resulting
continuous bijection on compact Hausdorff analytifications into a
homeomorphism. This is a topological fact, not PSH descent.

Pullback of genuine global PSH is unconditional. Projection formula and
BJ 7.2 give equal volume and energy,
\[E_{q^*\omega}(q^*\phi)=E_\omega(\phi).\] BJ 7.19 gives the actual
mixed operator identity \[q_*\bigwedge_i(q^*\omega_i+dd^cq^*\phi_i)
             =\bigwedge_i(\omega_i+dd^c\phi_i).                        \tag{25}\]
It follows by model projection formula and decreasing pairing
continuity, or by testing (3). Therefore a downstairs \(\mathcal E^1\)
potential pulls back to one upstairs. For a downstairs psh potential,
membership in \(\mathcal E^1\) is equivalent to that of its pullback by
the equal energy.

Actual onto PSH is a separate theorem with an actual antecedent: BJ 5.13
requires a sequence of larger classes \(\omega_m>\omega\), tending to
\(\omega\), each with EP. Its proof inserts small pulled-back
logarithmic ideal corrections, obtains downstairs potentials in those
larger classes, uses their increasing-net closure, then uses divisorial
determination. Under this hypothesis it is onto and a homeomorphism. The
equal energies then make it onto \(\mathcal E^1\), and (25) gives the
operator correspondence. In the unibranch homeomorphic case, probability
measures correspond bijectively; onto PSH gives equal suprema in (7), so
finite-dual-energy targets correspond as well. These onto statements
remain conditional here. Topological homeomorphism and the one-way
energy identities do not remove the EP hypothesis.

\subsection{What the countable-to-arbitrary-field passage would
prove}\label{what-the-countable-to-arbitrary-field-passage-would-prove}

Suppose the rational ample singular envelope theorem on countable
algebraically closed trivial subfields has actually been proved. Finite
geometry, line bundle, rational model obstacle, and model competitor
data descend to such a \(k_1\subset k\), since the algebraic closure of
the finitely generated coefficient field inside \(k\) is countable.
These trivial fields are complete. Reducedness, irreducibility and
ampleness descend faithfully. Over characteristic-zero perfect fields
finite-type normal algebras are geometrically normal. Their finite
normalization therefore commutes with these field extensions: its
extension is finite normal birational and hence the normalization. A
geometric second branch would persist after extension; thus geometric
unibranchness descends.

Let \(u_1=P_{\omega_1}(f_1)\) be the assumed continuous envelope for a
rational model-PL obstacle. Pullback \(u\) is a genuine global psh
minorant. The envelope of \(u_1+t(f_1-u_1)\) remains \(u_1\); continuous
genuine energy differentiation gives \(\int(f_1-u_1)\MA(u_1)=0\). The
measure is on contact. A genuine global model competitor below \(f\)
descends to a larger countable algebraically closed \(k_2\). Relative
nefness descends by the curve-degree test and faithful extension, and
ampleness by faithful flatness. Analytic restriction is surjective; BJ
LF755--758 supplies a continuous section for complete field extension of
the original trivial algebraically closed field. Thus the descended
competitor is below the descended obstacle.

The needed genuine mixed-measure field pushforward has an independent
model check. Regular characteristic-zero extension preserves vertical
\(R_1\) and \(S_2\): regular local fibers preserve regularity at
vertical generic points and the depth inequality at vertical points.
Hence BJ 1.13 preserves relative integral closure in the given generic
fiber. Central integral components over algebraically closed \(k_1\) are
geometrically integral. They do not split, their DVR orders of \(t\) are
unchanged, and their intersection degrees remain equal by flat extension
of cycles. Their divisorial restrictions are the original \(v_E\). Thus
(1) has the same pushed-forward weights. Decreasing model approximation
and testing continuous functions give
\[r_{21*}M_{\omega_2}^{\rm raw}(r_{21}^*u_1)
                           =M_{\omega_1}^{\rm raw}(u_1).               \tag{26}\]
This does not absolutely normalize the nonnormal generic fiber away.
Positivity and (26) put the upstairs measure on pulled-back contact, so
the global competitor is below the pulled-back envelope a.e. BJ global
domination 10.6 is valid on this singular irreducible \(k_2\)-variety
and gives the inequality everywhere. Model suprema, the envelope
contraction in sup norm, and PL density give continuous obstacles over
arbitrary \(k\). No singular BFJ Dirac equation is used.

For ample real classes, choose rational ample \(\omega_j\to\omega\) with
\(\omega_j-\omega\) ample. The genuine class closure and inclusions
\(\PSH(\omega)\subset\PSH(\omega_j)\) give EP in \(\omega\) once the
rational antecedent is proved. These are conditional passages; neither
produces the missing countable rational singular base. Applying BJ 12.8
and 5.13 after that base would settle the raw equation, all
normalization levels and arbitrary-support targets, but this paper has
not reached it.

\subsection{A nonsolvable finite-energy target on a nodal
curve}\label{a-nonsolvable-finite-energy-target-on-a-nodal-curve}

The obstruction outside unibranchness is genuine and unconditional. Let
\(C\) be an irreducible rational projective nodal curve,
\(q:\mathbb P^1\to C\) its normalization, and \(p_1\ne p_2\) the node
preimages. Put \(d=\deg L>0\) and \(v=\operatorname{ord}_{p_1}\) on the
common function field. It is a downstairs divisorial point and
\(\delta_v\) has finite dual energy by BJ 11.1.

Blow up \(\mathbb P^1\times\mathbb A^1\) at \((p_1,0)\) with exceptional
\(E\). The model class \(p^*q^*L+d\mathcal X_0-dE\) has degree zero on
the strict central component and degree \(d\) on \(E\), both
multiplicity one. These are all vertical curves, so it is relatively
nef; (1) gives raw measure \(d\delta_v\). Its bounded model potential is
\(g(s\operatorname{ord}_{p_1})=d\max(1-s,0)\) on the chosen ray and
\(d\) on the other rays. In particular its endpoint values at
\(p_1,p_2\) are \(0,d\).

If a downstairs \(\mathcal E^1\) solution to \(d\delta_v\) existed, its
pullback would have finite energy. Identity (25), positivity and total
mass force its upstairs measure to be \(d\delta_v\), since a valuation
on the common function field has a unique analytic lift. Upstairs
uniqueness on \(\mathbb P^1\) forces its potential to be \(g+c\). But a
pullback from \(C\) has equal values at the two node-preimage endpoints,
whereas \(g+c\) differs by \(d>0\). This is a contradiction, agreeing
with BJ 12.11. It is not a counterexample on a unibranch variety and
does not change the requested domain.

\section{Exact gaps, evidence, and rerunning the
deliverable}\label{exact-gaps-evidence-and-rerunning-the-deliverable}

The full original result remains unproved. The indispensable open gates
in this reconstruction are the following.

\begin{enumerate}
\def\labelenumi{\arabic{enumi}.}
\tightlist
\item
  For the whole corrected higher-dimensional relative local PL class,
  establish finite rational-domain representation and the required
  calibrated test-coordinate refinement. Formula (22) proves the
  finite-chart sign, not that every local input meets its analytic
  hypotheses. The strict/non-strict change in CLD 17.8.4 versus 18.8.8
  cannot be silently ignored. The smooth Laurent-field route requires
  its own strict-domain hypotheses.
\item
  Complete common genuine bounded-local representation, local mixed
  product agreement, weighted overlap independence, partition and
  integration-by-parts extensions, total mass and comparison for all
  smooth local inputs. The model-atom domination argument only uses
  these when supplied. Complete the corresponding smooth nef and weak
  affine base, and the controlled principal-open removable extension
  with the original exhaustion's upper growth.
\item
  Prove Rossi's arbitrary analytic-neighborhood application and split
  analytic-piece gluing in W 6.15, then the full local-projection,
  Hartogs/Runge and singular-stratum induction 6.3--6.5. The fixed curve
  trace equality and local \(\gamma\) repair do not construct those
  gluing functions or their smooth trace extensions.
\item
  Establish actual singular envelopes on the countable rational base.
  Only then do the conditional cardinality and real-class passages, BJ
  variational existence, and onto normalization furnish the unchanged
  original theorem. Arbitrary-support solvability and onto PSH/E1 are
  not asserted unconditionally here.
\end{enumerate}

The strongest projective-domain partial theorem is the genuine mixed
finite-energy operator, exact energy and raw mass, necessary
finite-dual-energy target condition, mixed integrability/continuity,
comparison, global domination and one-constant uniqueness, with the
exact EP equivalence for existence. The independently reconstructed
curve, trace, model-atom, support, collar and pole-order lemmas are new
retained proof inputs, with their stated limitations. No unrelated
valuation, characteristic, nonample, nonprojective or analytic equation
axis is inferred.

The preserved inputs are:

\begin{itemize}
\tightlist
\item
  \texttt{input-original.json}, byte-identical to \texttt{input.json};
\item
  \texttt{assignment-original.txt}, the original assignment value;
\item
  \texttt{task-original.json}, the original task value.
\end{itemize}

This proof is \texttt{P16.1-argument.md}; its standalone TeX rendering
is \texttt{P16.1-paper.tex}. \texttt{status.json} has exactly the
requested five keys and records \texttt{partial}.
\texttt{mathematical-audit.json} distinguishes unconditional source
deductions, independent proof drafts, conditional implications and
remaining gates.

\texttt{prepare\_provenance.py} verifies the twelve authorized PDF/text
hashes and saves exact versions, literal statements and relevant proof
bodies in \texttt{source-extracts/source-provenance.json}. Ranges are
inclusive one-based indices of UTF-8 text split by LF using
\texttt{split(\textquotesingle{}\textbackslash{}n\textquotesingle{})}.
All six source PDF/text pairs are retained in \texttt{sources/}. No
whole-source audit is claimed. \texttt{build.py} uses existing Pandoc
and runs existing pdfLaTeX twice, preserving stdout/stderr and separate
\texttt{tex-run-1.log} and \texttt{tex-run-2.log}. \texttt{verify.py}
checks preservation, hashes, LF excerpts, exact status schema,
earlier-only note support and build results. \texttt{environment.json}
preserves tool paths, Python/platform and working directory. Run
\texttt{python3\ rerun.py} in this directory.
\texttt{algebra\_checks.py} saves five exact rational-radius examples,
the two-by-two determinant and the length-two center in
\texttt{algebra-checks.json}; they are finite examples, not tests of an
infinite analytic theorem. Programs only do small text and document
operations; no heavy compute service is needed. No other campaign,
worker, model, software installation, credential, account/runtime
change, external message or external write was used.

Successful compilation, preservation and extraction are observations,
not mathematical proofs or novelty certificates. The notes and paper
remain unverified drafts. The final note does not claim a complete
solution of P16.1; \texttt{candidate=false} is required.

\section{References}\label{references}

{[}BJ{]} S. Boucksom and M. Jonsson, \emph{Global pluripotential theory
over a trivially valued field}, arXiv:1801.08229v4,
\url{https://arxiv.org/pdf/1801.08229v4}.

{[}W{]} L. Wu, \emph{On the local theory of non-archimedean psh
functions}, arXiv:2609.27610v1,
\url{https://arxiv.org/pdf/2609.27610v1}. Public v1; no independent
correctness verdict is assumed.

{[}CLD{]} A. Chambert-Loir and A. Ducros, \emph{Formes differentielles
reelles et courants sur les espaces de Berkovich}, arXiv:1204.6277v2,
supplied cache, \url{https://arxiv.org/pdf/1204.6277v2}.

{[}GM{]} W. Gubler and F. Martin, \emph{On Zhang's semipositive
metrics}, arXiv:1608.08030v3, \url{https://arxiv.org/pdf/1608.08030v3}.

{[}Wa{]} V. Wanner, \emph{Comparison of two notions of subharmonicity on
non-Archimedean curves}, arXiv:1801.04713v3,
\url{https://arxiv.org/pdf/1801.04713v3}.

{[}Th{]} A. Thuillier, \emph{Theorie du potentiel sur les courbes en
geometrie analytique non archimedienne. Applications a la theorie
d'Arakelov}, tel-00010990v1,
\url{https://theses.hal.science/tel-00010990v1/document}.

\end{document}
